%% =====================================================================
%%  Unit 19 — A trip to Boston  (Raconter un voyage)
%%  Pistes : 19.01–19.04 dialogues · 19.05–19.06 vocabulaire · 19.07–19.08 grammaire
%%           19.09 prononciation · 19.10–19.13 exercices (textes en annexe)
%% =====================================================================

\newcommand{\otIllusUnXIX}{\otIllusDuo{went/104/86/15/0.10,Boston/190/92/20/0.09,saw/122/14/13/0.08,
  train/150/8/11/0.08,bought/196/22/12/0.07,took/100/50/12/0.07}{jack}{How was Boston?}{lea}{It was amazing!}}

\unite{19}{A trip to Boston}{Raconter un voyage}{otIllusUnXIX}

\situation{Épisode 19}{Mi-février, un long week-end de trois jours. Léa et Emma ont pris le train pour
Boston, la plus européenne des villes américaines. Le mardi, au studio, Jack veut tout savoir… Mais
revenons d'abord au lundi soir, dans le train du retour.}

\begin{objectifs}
\item raconter un voyage au prétérit ;
\item employer les verbes irréguliers les plus courants (\en{went}, \en{saw}, \en{took}…) ;
\item poser des questions sur un voyage : \en{Where did you go?}, \en{How long did it take?} ;
\item parler de transports et d'hébergement ;
\item prononcer correctement les prétérits irréguliers.
\end{objectifs}

\auprogramme
  {Raconter un voyage, poser des questions sur le passé}
  {Le prétérit des verbes irréguliers}
  {Verbes irréguliers, voyage : 26 expressions}
  {Les verbes irréguliers les plus courants}
  {Boston et la Nouvelle-Angleterre}

\semaine{casque/otNavy/Dialogues/Écouter et lire les deux scènes,
  bulle/otGreenDark/Vocabulaire/Apprendre les 26 expressions,
  blocs/otViolet/Grammaire/Les verbes irréguliers,
  sons/otOrangeDark/Sons et culture/Les prétérits et Boston,
  crayon/otNavy/Exercices/12 exercices corrigés,
  coche/otGreenDark/Bilan/L'essentiel et le test express}

\clearpage
%% ---------------------------------------------------------------------
\rubrique{ecouter}{Un week-end à Boston}

\begin{dialogue}{1}{What did you buy?}{Lundi soir, dans le train qui ramène Léa et Emma de Boston à New
  York. Les sacs de souvenirs sont posés sur les sièges.}{19.01}{19.02}
\dit{emma}{What a weekend! So, Léa, \hl{what did you buy}?\note{1}}{Quel week-end ! Alors, Léa, qu'est-ce que tu as acheté ?}
\dit{lea}{I \hl{bought} a Red Sox cap for Jack and a mug for Sam.\note{2}}{J'ai acheté une casquette des Red Sox pour Jack et une tasse pour Sam.}
\dit{emma}{Ha! Jack is going to love it. And I \hl{took} two hundred photos!}{Ha ! Jack va adorer. Et moi, j'ai pris deux cents photos !}
\dit{lea}{I know! You \hl{took} a photo of every church!}{Je sais ! Tu as pris en photo chaque église !}
\dit{emma}{Well, I \hl{saw} the Freedom Trail, and I \hl{had} my camera… What did you prefer?}{Eh bien, j'ai vu la Freedom Trail, et j'avais mon appareil photo… Qu'est-ce que tu as préféré ?}
\dit{lea}{Harvard! We \hl{met} that nice student, and he \hl{gave} us a tour.\note{3}}{Harvard ! On a rencontré cet étudiant sympa, et il nous a fait visiter.}
\dit{emma}{Yes! And the clam chowder at Quincy Market… I \hl{ate} two bowls!}{Oui ! Et la soupe de palourdes au Quincy Market… J'en ai mangé deux bols !}
\dit{lea}{Wait… where's my ticket, Emma? \hl{I didn't put} it in my bag!\note{4}}{Attends… où est mon billet, Emma ? Je ne l'ai pas mis dans mon sac !}
\dit{emma}{Relax! It's on your phone. You \hl{got} an e-ticket!}{Détends-toi ! Il est sur ton téléphone. Tu as eu un billet électronique !}
\end{dialogue}

\begin{comprendre}
\pointnote{1} \en{\textbf{What did you buy?}} Les questions se forment toujours avec \en{did} + base
  verbale, même avec un verbe irrégulier.
\pointnote{2} \en{\textbf{I bought…}} \en{buy} est irrégulier : son prétérit est \en{bought}. Les verbes
  irréguliers ne prennent pas \en{-ed} : il faut apprendre leur forme (liste en annexe).
\pointnote{3} \en{\textbf{We met…, he gave us…}} \en{meet} → \en{met}, \en{give} → \en{gave}. Comme pour
  les verbes réguliers, la forme est la même pour toutes les personnes.
\pointnote{4} \en{\textbf{I didn't put it…}} À la forme négative, on revient à la base verbale : \en{I
  didn't put}, \en{she didn't buy}, \en{we didn't go}.
\end{comprendre}

\begin{dialogue}{2}{How was your trip?}{Mardi matin, au studio. Jack porte déjà sa casquette des Red
  Sox.}{19.03}{19.04}
\dit{jack}{Thanks for the cap, Léa! So, \hl{how was your trip}?}{Merci pour la casquette, Léa ! Alors, ce voyage ?}
\dit{lea}{It was amazing! We \hl{left} on Saturday morning and we \hl{took} the train.}{C'était génial ! On est parties samedi matin et on a pris le train.}
\dit{jack}{\hl{How long did it take}?\note{5}}{Ça a pris combien de temps ?}
\dit{lea}{About four hours. We \hl{got} there at noon.}{Environ quatre heures. On est arrivées à midi.}
\dit{jack}{\hl{Where did you stay}?}{Où est-ce que vous avez logé ?}
\dit{lea}{In a small hotel near the harbor. It wasn't expensive.}{Dans un petit hôtel près du port. Ce n'était pas cher.}
\dit{jack}{\hl{Did you go} to a Red Sox game?}{Vous êtes allées voir un match des Red Sox ?}
\dit{lea}{\hl{No, we didn't}. It's winter, Jack! But we \hl{went} to Harvard.\note{6}}{Non. C'est l'hiver, Jack ! Mais on est allées à Harvard.}
\end{dialogue}

\begin{comprendre}
\pointnote{5} \en{\textbf{How long did it take?}} « Ça a pris combien de temps ? » : \en{take} (prendre,
  durer) → \en{took}. Autres questions utiles : \en{How much did it cost?}, \en{Where did you go?}
\pointnote{6} \en{\textbf{We went to Harvard.}} \en{go} → \en{went} : la forme la plus irrégulière de
  toutes ! Mais à la négation et dans les questions : \en{We didn't go}, \en{Did you go?}
\end{comprendre}

%% ---------------------------------------------------------------------
\rubrique{mots}{Les mots de l'unité}

Voici les 26 expressions de l'unité. Les quatorze verbes du premier tableau sont parmi les plus
fréquents de l'anglais : apprenez-les par cœur, base verbale et prétérit ensemble. La liste complète
est en annexe, page~\pageref{annexe-Verbes}.

\begin{vocabulaire}{Verbes irréguliers}{19.05}
\categorie{Les indispensables}
\mot[go, went]{go · went}{ɡoʊ · went}{aller}
\mot[see, saw]{see · saw}{siː · sɑː}{voir}
\mot[take, took]{take · took}{teɪk · tʊk}{prendre}
\mot[eat, ate]{eat · ate}{iːt · eɪt}{manger}
\mot[have, had]{have · had}{hæv · hæd}{avoir}
\mot[buy, bought]{buy · bought}{baɪ · bɑːt}{acheter}
\mot[get, got]{get · got}{ɡet · ɡɑːt}{obtenir ; arriver}
\categorie{À connaître aussi}
\mot[come, came]{come · came}{kʌm · keɪm}{venir}
\mot[make, made]{make · made}{meɪk · meɪd}{faire, fabriquer}
\mot[meet, met]{meet · met}{miːt · met}{rencontrer}
\mot[leave, left]{leave · left}{liːv · left}{partir, quitter}
\mot[think, thought]{think · thought}{θɪŋk · θɑːt}{penser}
\mot[say, said]{say · said}{seɪ · sed}{dire}
\mot[give, gave]{give · gave}{ɡɪv · ɡeɪv}{donner}
\end{vocabulaire}

\begin{vocabulaire}{En voyage}{19.06}
\categorie{Transports}
\mot{a train station}{ə ˈtreɪn ˌsteɪʃən}{une gare}
\mot{a round-trip ticket}{ə ˌraʊnd trɪp ˈtɪkɪt}{un billet aller-retour}
\mot{a one-way ticket}{ə ˌwʌn weɪ ˈtɪkɪt}{un aller simple}
\mot{a reservation}{ə ˌrezɚˈveɪʃən}{une réservation}
\categorie{Sur place}
\mot{a hotel room}{ə hoʊˈtel ˌruːm}{une chambre d'hôtel}
\mot{check in}{ˌtʃek ˈɪn}{s'enregistrer (à l'hôtel)}
\mot{a tour}{ə ˈtʊr}{une visite guidée}
\mot{take photos}{ˌteɪk ˈfoʊtoʊz}{prendre des photos}
\mot{a souvenir}{ə ˌsuːvəˈnɪr}{un souvenir (objet)}
\mot{the harbor}{ðə ˈhɑːrbɚ}{le port}
\categorie{Raconter}
\mot{How was your trip?}{ˌhaʊ wəz jʊr ˈtrɪp}{C'était comment, ton voyage ?}
\mot{It was amazing!}{ɪt wəz əˈmeɪzɪŋ}{C'était génial !}
\end{vocabulaire}

%% ---------------------------------------------------------------------
\rubrique{grammaire}{Le prétérit des verbes irréguliers}

\otIntertitre[otViolet]{1 · Une forme à apprendre… seulement à l'affirmatif}
\begin{center}\otTableau
\begin{tabular}{@{}l@{\hspace{8mm}}>{\enOn}l<{\enOff}@{\hspace{8mm}}l@{}}
\otTH{Forme} & \otTH{Exemple} & \otTH{Français}\\\midrule
affirmation      & We \textbf{went} to Boston. · She \textbf{saw} Harvard.        & nous sommes allés…\\
négation         & We \textbf{didn't go} to Boston.                                & nous ne sommes pas allés…\\
question         & \textbf{Did} you \textbf{go} to Boston? · Where \textbf{did} you \textbf{go}? & tu es allé… ?\\
réponses courtes & Yes, we \textbf{did}. · No, she \textbf{didn't}.                & oui / non\\\bottomrule
\end{tabular}\end{center}

\begin{regle}[La bonne nouvelle]
La forme irrégulière ne sert qu'à l'affirmatif. Dans les négations et les questions, \en{did} fait tout le
travail et le verbe reprend sa base verbale : \en{I bought} → \en{I didn't buy} → \en{Did you buy?}
\end{regle}

\blocaudio{19.07}{Écoutez et répétez}{Six phrases avec un verbe irrégulier.}

\otIntertitre[otViolet]{2 · Questionner sur un voyage}
\begin{center}\otTableau
\begin{tabular}{@{}>{\enOn}l<{\enOff}@{\hspace{6mm}}l@{\hspace{8mm}}>{\enOn}l<{\enOff}@{\hspace{6mm}}l@{}}
\otTH{Question} & \otTH{Français} & \otTH{Question} & \otTH{Français}\\\midrule
Where did you go?        & où es-tu allé ?           & How did you travel?      & comment as-tu voyagé ?\\
How long did it take?    & combien de temps ? & Where did you stay?   & où as-tu logé ?\\
How much did it cost?    & combien ça a coûté ?      & What did you see?        & qu'est-ce que tu as vu ?\\\bottomrule
\end{tabular}\end{center}

\blocaudio{19.08}{Écoutez et répétez}{Une question sur un voyage, puis sa réponse.}

\begin{attention}[Les pièges des francophones]
\fauxjuste{I goed to Boston.}{I went to Boston.}
\fauxjuste{Did you went to Harvard?}{Did you go to Harvard?}
\fauxjuste{I didn't saw the game.}{I didn't see the game.}
\fauxjuste{We buyed souvenirs.}{We bought souvenirs.}
\fauxjuste{Where you went?}{Where did you go?}
\end{attention}

%% ---------------------------------------------------------------------
\rubrique{sons}{Les verbes irréguliers les plus courants}

Les prétérits irréguliers ont souvent une prononciation inattendue. \en{bought}, \en{thought} et
\en{taught} riment ensemble \phon{ɑːt} : le \en{gh} ne se prononce jamais. \en{said} se dit \phon{sed},
et non « séd ». \en{ate} se prononce exactement comme \en{eight} \phon{eɪt}. Entraînez-vous avec ces
phrases jusqu'à ce que ces formes deviennent automatiques.

\begin{phrases}{19.09}{Écoutez chaque phrase et répétez-la. Le prétérit est en couleur ; sa prononciation
est indiquée à droite.}
\phrase[I bought a cap.]{I \textcolor{otOrangeDark}{bought} a cap.}{\en{buy} → \en{bought} \phon{bɑːt}}
\phrase[I thought about it.]{I \textcolor{otOrangeDark}{thought} about it.}{\en{think} → \en{thought} \phon{θɑːt}}
\phrase[She taught me English.]{She \textcolor{otOrangeDark}{taught} me English.}{\en{teach} → \en{taught} \phon{tɑːt}}
\phrase[He said hello.]{He \textcolor{otOrangeDark}{said} hello.}{\en{say} → \en{said} \phon{sed}}
\phrase[We saw the harbor.]{We \textcolor{otOrangeDark}{saw} the harbor.}{\en{see} → \en{saw} \phon{sɑː}}
\phrase[I ate clam chowder.]{I \textcolor{otOrangeDark}{ate} clam chowder.}{\en{eat} → \en{ate} \phon{eɪt}, comme \en{eight}}
\end{phrases}

%% ---------------------------------------------------------------------
\rubrique{culture}{Boston et la Nouvelle-Angleterre}

\begin{culture}[Le berceau de l'Amérique]
Fondée en 1630, Boston est l'une des plus anciennes villes des États-Unis. C'est là qu'a commencé la
révolution américaine. Pour la découvrir, suivez la \en{Freedom Trail} : une ligne rouge tracée sur
les trottoirs, longue de 4\,km (2,5 miles), relie seize sites historiques, du \en{Boston Common} (1634),
le plus ancien parc public du pays, à l'\en{Old North Church}.
\end{culture}

\begin{paire}
\begin{culture}[Harvard et le MIT]
De l'autre côté de la rivière Charles, à Cambridge, se trouvent Harvard, la plus ancienne université
des États-Unis (1636), et le MIT, célèbre pour ses ingénieurs. Les visites du campus sont souvent
guidées par des étudiants.
\end{culture}
\begin{culture}[Le Boston des gourmands et des sportifs]
Spécialités : la \en{clam chowder} (soupe crémeuse aux palourdes) et le homard (\en{lobster}). Les fans
de baseball se pressent à Fenway Park (1912), le plus ancien stade de la ligue, où jouent les Red Sox.
\end{culture}
\end{paire}

%% ---------------------------------------------------------------------
\rubrique{pratique}{Exercices}

\exercice[ecrire,niveau=1]{Les prétérits irréguliers}{Écrivez le prétérit de chaque verbe.}
\begin{anglais}
\begin{questions}[4]
\item go → \trou[18mm]{went}
\item see → \trou[18mm]{saw}
\item eat → \trou[18mm]{ate}
\item buy → \trou[18mm]{bought}
\item take → \trou[18mm]{took}
\item have → \trou[18mm]{had}
\item meet → \trou[18mm]{met}
\item think → \trou[18mm]{thought}
\end{questions}
\end{anglais}

\exercice[ecrire,niveau=2]{Le message de Léa}{Léa écrit à ses parents. Complétez avec le prétérit des
verbes entre parenthèses.}
\begin{anglais}
\begin{tcolorbox}[otBase,colback=otNavyTint,colframe=otNavyTint,fontupper=\fontsize{11}{19}\selectfont]
Dear Mom and Dad,\par
Last weekend, Emma and I \trou[20mm]{went} \textit{(go)} to Boston. We \trou[20mm]{took} \textit{(take)}
the train on Saturday. We \trou[20mm]{saw} \textit{(see)} the old churches and the harbor. I
\trou[20mm]{ate} \textit{(eat)} clam chowder, a delicious soup, and I \trou[20mm]{bought} \textit{(buy)}
a cap for Jack. We \trou[20mm]{had} \textit{(have)} a wonderful time! Love, Léa
\end{tcolorbox}
\end{anglais}

\exercice[ecrire,niveau=2]{Négations et questions}{Transformez chaque phrase : (?) en question, (−) en
phrase négative.}
\begin{questions}
\item We went to Harvard. (−) \hfill \enb{\trou[68mm]{We didn't go to Harvard.}}
\item She bought a cap. (?) \hfill \enb{\trou[68mm]{Did she buy a cap?}}
\item They saw the game. (−) \hfill \enb{\trou[68mm]{They didn't see the game.}}
\item You took the train. (?) \hfill \enb{\trou[68mm]{Did you take the train?}}
\item He ate lobster. (−) \hfill \enb{\trou[68mm]{He didn't eat lobster.}}
\item Emma met a student. (?) \hfill \enb{\trou[68mm]{Did Emma meet a student?}}
\end{questions}

\exercice[ecrire,niveau=2]{Posez la question}{Voici des réponses. Écrivez la question qui porte sur les
mots soulignés.}
\begin{questions}
\item We went to \underLine{Boston}. \hfill \enb{\trou[66mm]{Where did you go?}}
\item We stayed \underLine{three days}. \hfill \enb{\trou[66mm]{How long did you stay?}}
\item It cost \underLine{120 dollars}. \hfill \enb{\trou[66mm]{How much did it cost?}}
\item We traveled \underLine{by train}. \hfill \enb{\trou[66mm]{How did you travel?}}
\item We saw \underLine{the Freedom Trail}. \hfill \enb{\trou[66mm]{What did you see?}}
\end{questions}

\exercice[lire,niveau=2]{Le Boston Tea Party}{Lisez le texte, puis indiquez si chaque phrase est vraie
(V) ou fausse (F).}
\begin{anglais}
\begin{tcolorbox}[otBase,colback=white,colframe=otRule,boxrule=0.5pt,fontupper=\fontsize{10.8}{14.4}\selectfont]
{\otSemi\fontsize{11}{13}\selectfont The Boston Tea Party}\par\vspace{1.4mm}
In 1773, Boston was an English colony. The colonists didn't want to pay a new tax on tea. On the night
of December 16th, a group of men went onto three English ships in the harbor. They threw 342 boxes of
tea into the water! The English king was very angry. It was the beginning of the American Revolution.
\par\vspace{1mm}{\fontsize{8.6}{10}\selectfont\color{otMuted}\foreignlanguage{french}{\textit{a tax} : un impôt, une taxe · \textit{a ship} : un navire · \textit{threw} : prétérit de \textit{throw}, jeter · \textit{angry} : en colère · \textit{the beginning} : le début}}
\end{tcolorbox}
\begin{questions}[2]
\item In 1773, Boston was an English colony. \hfill\trou[9mm]{V}
\item The colonists wanted to pay the tax. \hfill\trou[9mm]{F}
\item The men went onto three ships. \hfill\trou[9mm]{V}
\item They drank the tea. \hfill\trou[9mm]{F}
\item It was at night. \hfill\trou[9mm]{V}
\item It was the end of the Revolution. \hfill\trou[9mm]{F}
\end{questions}
\end{anglais}

\exercice[lire,niveau=1]{Questions et réponses}{Associez chaque question à sa réponse : écrivez la
lettre.}
\begin{anglais}
\noindent\begin{minipage}[t]{0.5\linewidth}
\begin{questions}[1][english]
\item Where did you go? \hfill\trou[9mm]{c}
\item How did you travel? \hfill\trou[9mm]{a}
\item Where did you stay? \hfill\trou[9mm]{e}
\item What did you see? \hfill\trou[9mm]{b}
\item How long did you stay? \hfill\trou[9mm]{f}
\item Did you like it? \hfill\trou[9mm]{d}
\end{questions}
\end{minipage}\hfill
\begin{minipage}[t]{0.42\linewidth}
\begin{questions}[1][english]
\item[a.] By train.
\item[b.] The Freedom Trail.
\item[c.] To Boston.
\item[d.] Yes, I did!
\item[e.] In a small hotel.
\item[f.] Three days.
\end{questions}
\end{minipage}
\end{anglais}

\exercice[ecouter,niveau=1,piste=19.10]{Présent ou passé ?}{Écoutez chaque phrase, dite deux fois.
Est-elle au présent ou au passé ? Cochez.}
\begin{questions}[3]
\item \choix{présent}\quad\choix{passé}
\item \choix{présent}\quad\choix{passé}
\item \choix{présent}\quad\choix{passé}
\item \choix{présent}\quad\choix{passé}
\item \choix{présent}\quad\choix{passé}
\item \choix{présent}\quad\choix{passé}
\end{questions}

\exercice[ecouter,ecrire,niveau=2,piste=19.11]{Le récit de Léa}{Mrs.\ Lopez demande à Léa comment
s'est passé son voyage. Écoutez et complétez.}
\begin{anglais}
\begin{questions}
\item Léa and Emma left on \trou[30mm]{Saturday}.
\item They took the \trou[30mm]{train}.
\item The trip took \trou[30mm]{four} hours.
\item They stayed in a small \trou[30mm]{hotel}.
\item Léa ate clam chowder and \trou[30mm]{lobster}.
\end{questions}
\end{anglais}

\exercice[ecouter,ecrire,niveau=2,piste=19.12]{Dictée à trous}{Sam et Léa parlent du voyage. Écoutez et
complétez.}
\begin{anglais}
\noindent\begin{tabular}{@{}l@{\hspace{6mm}}l@{}}
\otNomDit{sam} & Thanks for the mug! Where did you \trou{buy} it?\\[1.8mm]
\otNomDit{lea} & I \trou{bought} it in a small shop near the harbor.\\[1.8mm]
\otNomDit{sam} & \trou{Did} you see the Old North Church?\\[1.8mm]
\otNomDit{lea} & Yes, we \trou{did}. Emma \trou{took} a hundred photos!\\[1.8mm]
\otNomDit{sam} & Ha! And what did you \trou{eat}?\\[1.8mm]
\otNomDit{lea} & We \trou{ate} clam chowder. We \trou{had} a great time!\\
\end{tabular}
\end{anglais}

\exercice[ecrire,niveau=2]{Traduisez}{Traduisez en anglais.}
\begin{questions}[1][french]
\item Nous sommes allés à Boston le week-end dernier.\lignes{1}
\item J'ai pris le train et j'ai vu la mer.\lignes{1}
\item Tu as acheté des souvenirs ? — Oui. \textit{(réponse courte)}\lignes{1}
\item Elle n'a pas mangé de homard.\lignes{1}
\item Ça a pris combien de temps ?\lignes{1}
\item Où avez-vous logé ?\lignes{1}
\item Nous avons rencontré un étudiant sympa.\lignes{1}
\item C'était génial !\lignes{1}
\end{questions}

\exercice[ecrire,niveau=3]{Mon dernier voyage}{En six phrases, racontez votre dernier voyage : où vous
êtes allé, comment, où vous avez logé, ce que vous avez vu et mangé.}
\lignes{6}

\exercice[ecouter,parler,niveau=3,piste=19.13]{À vous de parler !}{Jack vous interroge sur votre
dernier voyage. Répondez \textbf{à voix haute} pendant chaque pause.}
\begin{center}\otTableau
\begin{tabular}{@{}>{\otSemi\fontsize{8.6}{10}\selectfont\color{otNavy}}l@{\hspace{3mm}}>{\enOn}p{70mm}<{\enOff}@{\hspace{6mm}}p{78mm}@{}}
 & \otTH{Jack dit\ldots} & \otTH{Vous\ldots}\\\midrule
1. & Where did you go on your last trip? & dites où.\\
2. & How did you travel? & dites comment.\\
3. & Where did you stay? & dites où vous avez logé.\\
4. & What did you see? & citez une ou deux visites.\\
5. & Did you buy any souvenirs? & répondez et précisez.\\\bottomrule
\end{tabular}\end{center}

%% ---------------------------------------------------------------------
\rubrique{bilan}{Faites le point}

\begin{essentiel}
\begin{minipage}[t]{0.54\linewidth}
\otEssTitre{Les phrases clés}
\renewcommand{\arraystretch}{1.32}
\begin{tabular}{@{}>{\enOn\otSemi\fontsize{9.4}{11}\selectfont}l<{\enOff}@{\hspace{3mm}}>{\fontsize{9.2}{11}\selectfont\color{white!78!otNavy}}l@{}}
How was your trip?         & C'était comment, ton voyage ?\\
We went to Boston.         & Nous sommes allés à Boston.\\
We took the train.         & Nous avons pris le train.\\
Did you see Harvard?       & Tu as vu Harvard ?\\
I didn't buy souvenirs.    & Je n'ai pas acheté de souvenirs.\\
It was amazing!            & C'était génial !\\
\end{tabular}
\end{minipage}\hfill
\begin{minipage}[t]{0.41\linewidth}
\otEssTitre{Les irréguliers}
{\fontsize{9.2}{11.6}\selectfont \en{went}, \en{saw}, \en{took}, \en{ate}, \en{had}, \en{bought},
\en{got}, \en{met}, \en{left}… (liste en annexe)\par}
\otEssTitre{Négation et question}
{\fontsize{9.2}{11.6}\selectfont \en{didn't} / \en{Did…?} + base verbale : \en{I didn't go} · \en{Did
you go?}\par}
\otEssTitre{Prononciation}
{\fontsize{9.2}{11.6}\selectfont \en{bought}, \en{thought} \phon{ɑːt} · \en{said} \phon{sed} ·
\en{ate} = \en{eight}\par}
\end{minipage}
\end{essentiel}

\begin{jesais}
\item raconter un voyage au prétérit ;
\item employer une quinzaine de verbes irréguliers courants ;
\item poser des questions sur un voyage avec \en{did} ;
\item parler de transports et d'hébergement ;
\item prononcer \en{bought}, \en{said}, \en{saw} sans hésiter ;
\item présenter Boston et la \en{Freedom Trail}.
\end{jesais}

\clearpage
\begin{testexpress}
\qcm[c]{Le prétérit de \en{go} est :}{\en{goed}}{\en{gone}}{\en{went}}
\qcm[a]{« Nous avons pris le train » se dit :}{\en{We took the train.}}{\en{We taked the train.}}{\en{We take the train.}}
\qcm[b]{« Tu es allé à Harvard ? » se dit :}{\en{Did you went to Harvard?}}{\en{Did you go to Harvard?}}{\en{You went to Harvard?}}
\qcm[c]{« Je n'ai pas vu le match » se dit :}{\en{I didn't saw the game.}}{\en{I not saw the game.}}{\en{I didn't see the game.}}
\qcm[a]{Le prétérit de \en{buy} est :}{\en{bought}}{\en{buyed}}{\en{brought}}
\qcm[b]{« Ça a coûté combien ? » se dit :}{\en{How much it cost?}}{\en{How much did it cost?}}{\en{How many did it cost?}}
\qcm[a]{\en{said} se prononce :}{\phon{sed}}{\phon{seɪd}}{\phon{saɪd}}
\qcm[c]{\en{ate} se prononce comme :}{\en{at}}{\en{eat}}{\en{eight}}
\qcm[b]{Un billet aller-retour se dit :}{\en{a one-way ticket}}{\en{a round-trip ticket}}{\en{a back ticket}}
\qcm[c]{La \en{Freedom Trail} est :}{un train historique}{une université}{un parcours historique à Boston}
\end{testexpress}

\noindent{\fontsize{9.6}{12}\selectfont\color{otMuted}Niveau visé (CECRL, niveau A1) : « peut raconter
un voyage en phrases simples ». Les corrigés de l'unité sont page~\pageref{cor-u19-1} ; les textes des
pistes d'exercices, page~\pageref{trans-19.10}.\par}
